\documentclass{situationist-preprint}

\title{Survival Without Memory}
\msid{PID-08-a}
\msfamily{PID-08}
\msdiscipline{Personal Identity and the Self}
\mstype{Principal manuscript}
\mskeywords{amnesia; memory criterion; psychological continuity; Locke; Reid; Butler; character}

\begin{document}
\maketitle

\begin{abstract}
Can a person survive total amnesia, the loss of all autobiographical memory? We argue that she can. Survival does not require memory continuity, only the continuity of other psychological capacities: character, values, skills, semantic knowledge, emotional dispositions and the capacity for self-conscious thought. The memory criterion descending from Locke fails for familiar reasons, transitivity and circularity, and the psychological criterion of PID-01, properly stated, does not include memory as a necessary component. Clinical cases of profound amnesia show that the other capacities can continue when memory does not, and the judgments these cases elicit, from clinicians, families and the patients themselves, treat the patient as the same person. We state the criterion, derive the conclusion and reply to the objections that without memory there is no self-concern, no narrative and no criterion that excludes duplicates.
\end{abstract}

\section{The problem}

\citet[II.xxvii.9]{Locke1694} held that as far as consciousness can be extended backwards to any past action or thought, so far reaches the identity of that person. Read as a memory criterion, the view faced two objections at once. \citet{Butler1736} argued that memory presupposes identity and so cannot constitute it, and \citet[III.6]{Reid1785} that the criterion violates transitivity: the old general remembers the young officer's deeds, the officer the boy's flogging, but the general does not remember the flogging. Neo-Lockean views answered both objections by moving from memory to psychological continuity, the ancestral of overlapping connections, and by replacing memory with quasi-memory \citep{Shoemaker1970,Parfit1984}. But they generally kept memory as the central kind of connection.

Clinical amnesia puts the question directly. Clive Wearing, after herpes encephalitis in 1985, lost almost all autobiographical memory and the ability to form new episodic memories, while keeping his musical skills, his character and his devotion to his wife \citep{Wearing2005}. The patient H.M. could form no new long-term declarative memories after surgery in 1953 \citep{ScovilleMilner1957}. If memory were necessary, such patients would not survive their illness.

\section{Framework}\label{s:frame}

\begin{definition}\label{d:amnesia}
\emph{Total amnesia} is the loss of all episodic and autobiographical memory, retrograde and anterograde, with the preservation of other psychological capacities.
\end{definition}

Total amnesia so defined is a limiting case; actual patients retain fragments. The distinction it presupposes, between episodic memory of one's own past and the semantic and procedural systems, is standard in cognitive psychology \citep{Tulving1985}.

\begin{definition}\label{d:capacities}
The \emph{non-mnemonic psychological capacities} of a person are her character traits, values, desires, skills, semantic knowledge, emotional dispositions and capacity for self-conscious thought.
\end{definition}

\begin{assumption}\label{a:pc}
A later person is the same person as an earlier one if they are psychologically continuous by a non-branching process with the right kind of cause (PID-01), where psychological continuity is the ancestral of strong direct connections of any of the kinds listed in Definition~\ref{d:capacities} or of memory.
\end{assumption}

\begin{assumption}\label{a:strong}
The continuity of the non-mnemonic capacities can by itself constitute strong connectedness, when these capacities are rich and idiosyncratic and are causally sustained by the same realizer.
\end{assumption}

Assumption~\ref{a:pc} is the psychological criterion as \citet{Parfit1984} stated it, with direct connections of several kinds; memory is one kind among them, not a necessary condition. Assumption~\ref{a:strong} is supported by the specificity of character and value: a person's particular configuration of loves, tastes, convictions, fears and skills is as individuating as her memories, and in amnesia it is sustained by the same brain that sustained it before.

\section{The argument}\label{s:arg}

\begin{theorem}\label{t:main}
A person can survive total amnesia.
\end{theorem}

\begin{proof}
In total amnesia, by Definition~\ref{d:amnesia}, the non-mnemonic capacities are preserved and are sustained by the same brain. By Assumption~\ref{a:strong} their continuity constitutes strong connectedness between the person before and the person after; the process does not branch; so by Assumption~\ref{a:pc} the later person is the earlier one.
\end{proof}

\begin{corollary}\label{c:notnec}
Memory continuity is not necessary for personal survival.
\end{corollary}

The theorem fits the evidence about how people judge such cases. Families of amnesic patients regard them as the same person, changed; and experimental work finds that people judge identity to be more disrupted by changes in moral character than by loss of memory \citep{StrohmingerNichols2014}. The judgments are evidence, not proof, but they converge with the criterion.

The theorem is also what PID-05 predicts for the organism: the animal that you are survives amnesia, since its life continues. The present result concerns the person in the sense of PID-01, and it shows that the two verdicts agree in this case. By PID-06, connectedness carries most of what matters; total amnesia removes one strand of connectedness and leaves the others, so that the amnesic person keeps a substantial part of what matters in survival, not merely identity.

\section{Objections and replies}\label{s:obj}

\begin{objection}[No self-concern]
Without memories a person cannot recognize her past as hers, and so has nothing that makes her the same self.
\end{objection}

\begin{reply}
Recognition of one's past is a mark of a good life of a certain kind, not a condition of persistence. The amnesic person still loves what she loved, judges as she judged and acts from the same dispositions; that she cannot recall the occasions on which she acquired them does not make them someone else's.
\end{reply}

\begin{objection}[No narrative]
Personal identity is narrative identity, and without autobiographical memory there is no narrative \citep{Schechtman1996}.
\end{objection}

\begin{reply}
PID-09 concludes that narrative identity is real but distinct from numerical identity. The amnesic person loses much of her own capacity to tell her story, and others, who share the authorship of narratives according to PID-09, continue to tell it. What amnesia damages is narrative, which is a real loss; what survives is numerical identity.
\end{reply}

\begin{objection}[Duplicates]
Without memory, the criterion cannot distinguish the person from anyone with the same character and skills.
\end{objection}

\begin{reply}
Assumption~\ref{a:pc} requires continuity with the right kind of cause, and Assumption~\ref{a:strong} requires that the capacities be causally sustained by the same realizer. A stranger with a similar character is not causally continuous with the person; the amnesic person's character is the same character, sustained by the same brain.
\end{reply}

\section{What the result does not establish}\label{s:not}

The result does not show that a person survives the loss of all psychological capacities. If, besides memory, character, skills and the capacity for self-conscious thought were lost, as in a persistent vegetative state, nothing would sustain psychological continuity, and the person in the sense of PID-01 would not survive, though by PID-05 the organism would. Nor does the result say that amnesia leaves what matters intact; it removes a significant part of it.

\section{Conclusion}

Total amnesia removes memory continuity but leaves the other psychological capacities, and their continuity, sustained by the same brain, is enough for psychological continuity in the sense of PID-01. A person can survive total amnesia; memory is not necessary for survival.

\begin{thebibliography}{99}\small
\bibitem[Butler(1736)]{Butler1736} Butler, J. (1736). Of personal identity. First dissertation appended to \emph{The Analogy of Religion}.
\bibitem[Locke(1694)]{Locke1694} Locke, J. (1694). \emph{An Essay Concerning Human Understanding}, 2nd edn. Cited by book, chapter and section.
\bibitem[Parfit(1984)]{Parfit1984} Parfit, D. (1984). \emph{Reasons and Persons}. Oxford University Press.
\bibitem[Reid(1785)]{Reid1785} Reid, T. (1785). \emph{Essays on the Intellectual Powers of Man}. Cited by essay and chapter.
\bibitem[Schechtman(1996)]{Schechtman1996} Schechtman, M. (1996). \emph{The Constitution of Selves}. Cornell University Press.
\bibitem[Scoville and Milner(1957)]{ScovilleMilner1957} Scoville, W.~B., and Milner, B. (1957). Loss of recent memory after bilateral hippocampal lesions. \emph{Journal of Neurology, Neurosurgery and Psychiatry}, 20(1), 11--21.
\bibitem[Shoemaker(1970)]{Shoemaker1970} Shoemaker, S. (1970). Persons and their pasts. \emph{American Philosophical Quarterly}, 7(4), 269--285.
\bibitem[Strohminger and Nichols(2014)]{StrohmingerNichols2014} Strohminger, N., and Nichols, S. (2014). The essential moral self. \emph{Cognition}, 131(1), 159--171.
\bibitem[Tulving(1985)]{Tulving1985} Tulving, E. (1985). Memory and consciousness. \emph{Canadian Psychology}, 26(1), 1--12.
\bibitem[Wearing(2005)]{Wearing2005} Wearing, D. (2005). \emph{Forever Today: A Memoir of Love and Amnesia}. Doubleday.
\end{thebibliography}

\end{document}
